% ============================== PROBLEM 2 ============================
\settab{PROBLEM 2}{3}
\renewcommand{\chapribbontext}{Problem 2 \;·\; 7 pts}
\chapter[Problem 2: A Function Defined by Cases]{A Function\\ Defined by Cases}\label{ch:p2}

\begin{examnote}[title={The mock}]
$g(x,y)=\dfrac{xy^2}{x^2+y^2}$ for $(x,y)\ne(0,0)$, and $g(0,0)=0$.
\end{examnote}

\begin{idea}
This problem has \textbf{one golden rule}: away from the origin, use the normal rules; \textbf{at the
origin, never plug into a formula, always compute a limit}. Every question is written in two
paragraphs: ``For $(x,y)\ne(0,0)$'' and ``At $(0,0)$''.
\end{idea}

\section{The polar trick: know the answers in 30 seconds}

On scrap paper, put $x=r\cos\theta$, $y=r\sin\theta$:
\[
  g=\frac{r\cos\theta\cdot r^2\sin^2\theta}{r^2}=r\cdot\cos\theta\sin^2\theta .
\]
\noindent\begin{minipage}[c]{0.53\linewidth}
{\small
\renewcommand{\arraystretch}{1.3}
\begin{tabular}{@{}ll@{}}\toprule
$g=r^n\cdot(\text{angle})$ & at the origin\\ \midrule
$n=0$, e.g. $\frac{xy}{x^2+y^2}$ & not continuous\\[2pt]
$n=1$, e.g. the mock & continuous, \textbf{not diff.}\\[2pt]
$n\ge2$, e.g. $\frac{x^2y^2}{x^2+y^2}$ & differentiable (diff.)\\ \bottomrule
\end{tabular}
}
\par\medskip This is not a proof, so do not write it as your answer. It tells you \emph{which} proof sentence to use.
\end{minipage}\hfill
\begin{minipage}[c]{0.44\linewidth}\centering

\begin{tikzpicture}
\begin{axis}[bp3d, width=6.4cm, height=5.2cm, view={-40}{26}, xlabel=$x$, ylabel=$y$, zlabel=$z$,
   xmin=-1.25, xmax=1.25, ymin=-1.25, ymax=1.25, zmin=-0.6, zmax=0.6]
  \addplot3[gsurf, shader=interp, domain=0:1, y domain=0:360, samples=14, samples y=56,
     point meta=z] ({x*cos(y)},{x*sin(y)},{x*cos(y)*sin(y)^2});
  \foreach \t in {0,35,70,110,145,180,215,250,290,325}{
    \edef\tmp{\noexpand\addplot3[coralD, line width=0.9pt, domain=0:1, samples=2, samples y=1]
      ({x*cos(\t)},{x*sin(\t)},{x*cos(\t)*sin(\t)^2});}\tmp}
  \node[dot] at (axis cs:0,0,0) {};
\end{axis}
\end{tikzpicture}
\figcap{The graph of $g$ is made of straight lines through the origin (coral), each with its own
slope: a wavy cone. A cone tip has a slope in every direction but no tangent plane.}

\end{minipage}


% -------------------------------------------------------------------
\needspace{12\baselineskip}
\section{Q1. The partial derivatives}



\begin{writethis}
For $(x,y)\ne(0,0)$, quotient rule $\big(\frac uv\big)'=\frac{u'v-uv'}{v^2}$:
\[
 g_x=\frac{y^2(x^2+y^2)-xy^2\cdot2x}{(x^2+y^2)^2}=\frac{y^2(y^2-x^2)}{(x^2+y^2)^2},
\]
\[
 g_y=\frac{2xy(x^2+y^2)-xy^2\cdot2y}{(x^2+y^2)^2}=\frac{2x^3y}{(x^2+y^2)^2}.
\]
At $(0,0)$: $g(h,0)=0$ and $g(0,k)=0$, so
$g_x(0,0)=\lim\limits_{h\to0}\frac{0-0}{h}=0$ and $g_y(0,0)=\lim\limits_{k\to0}\frac{0-0}{k}=0$.
\end{writethis}



% -------------------------------------------------------------------
\section{Q2. Continuity}

\begin{idea}
Squeeze $|g|$ below something that goes to $0$, using $\frac{y^2}{x^2+y^2}\le1$. (Checking lines $y=mx$ is not a proof.)
\end{idea}

\begin{writethis}
For $(x,y)\ne(0,0)$, $g$ is a quotient of polynomials with nonzero denominator: continuous.

At $(0,0)$ \seeproof{4}: \ $|g(x,y)-0|=|x|\cdot\dfrac{y^2}{x^2+y^2}\le|x|\to0$, so
$\lim\limits_{(x,y)\to(0,0)}g=0=g(0,0)$. Hence $g$ is continuous on $\R^2$.
\end{writethis}



% -------------------------------------------------------------------
\section{Q3. Directional derivatives at the origin}

\begin{idea}
Slope at $O$ in direction $v=(a,b)$: the limit of $\frac{g(ta,tb)-g(0,0)}{t}$. Not $\grad g\cdot v$ (only valid if differentiable).
\end{idea}

\begin{writethis}
Let $a^2+b^2=1$. For $t\ne0$: \ $g(ta,tb)=\dfrac{t^3ab^2}{t^2(a^2+b^2)}=t\,ab^2$, so
$\dfrac{g(ta,tb)-0}{t}=ab^2$. The limit exists: \ $\boxed{D_vg(0,0)=ab^2}$.
\end{writethis}

\begin{watch}\watchout The tempting wrong answer is $\grad g(0,0)\cdot v=0$. The true answer $ab^2$ is not $0$ (try $a=b=\frac1{\sqrt2}$).\end{watch}

% -------------------------------------------------------------------
\section{Q4. Not differentiable at the origin}

\begin{idea}
Differentiable means: near $O$, $g$ is a plane plus an error that is tiny even after dividing by the
distance. Here the plane is $z=0$, so test $\frac{g(h,k)}{\text{distance}}\to0$. On the diagonal $h=k$ it stays constant: it fails.
(Never test on the axes: there it is $0$, which proves nothing.)
\end{idea}

\begin{writethis}
\seeproof{6} \ $g(0,0)=0$ and $\grad g(0,0)=(0,0)$. If $g$ were differentiable at $O$, then
$\varepsilon(h,k)=\dfrac{g(h,k)}{\sqrt{h^2+k^2}}=\dfrac{hk^2}{(h^2+k^2)^{3/2}}\to0$.
But for $h=k=t>0$: \ $\varepsilon(t,t)=\dfrac{t^3}{(2t^2)^{3/2}}=\dfrac{1}{2\sqrt2}\ne0$.
So $g$ is \textbf{not differentiable} at the origin.
\end{writethis}



% -------------------------------------------------------------------
\section{Q5. Is $g_x$ continuous? Is $g$ of class $C^1$?}

\begin{writethis}
\seeproof{7} \ On the $y$-axis ($x=0$, $y\ne0$): \ $g_x(0,y)=\dfrac{y^2\cdot y^2}{y^4}=1$, which does not tend to
$g_x(0,0)=0$. So $g_x$ is not continuous at the origin, and $g\notin C^1(\R^2)$.
(Consistent with Q4: $C^1$ would have forced differentiability.)

\smallskip\textbf{The whole story in one sentence:} $g$ is continuous with a slope in every direction at $O$,
but not differentiable there, so not $C^1$ either.
\end{writethis}


